% !PRINTED 342
\weberpage{342}
If $\alpha, \beta, \gamma \ldots$ are arbitrary real or complex quantities, subject however to the condition that the absolute value of $\alpha$ is \gesperrt{greater} than the absolute value of all the others, so that the absolute values of the fractions $\beta:\alpha, \gamma:\alpha \ldots$ are proper fractions, then
\weq{1}{
\frac{\alpha^m + \beta^m + \gamma^m + \cdots}{\alpha^{m-1}+\beta^{m-1}+\gamma^{m-1}+\cdots}
= \alpha\,\frac{1+\left(\dfrac{\beta}{\alpha}\right)^m+\left(\dfrac{\gamma}{\alpha}\right)^m+\cdots}{1+\left(\dfrac{\beta}{\alpha}\right)^{m-1}+\left(\dfrac{\gamma}{\alpha}\right)^{m-1}+\cdots},
}
and the greater $m$ becomes, the more this expression will approach the limit $\alpha$, as the second representation shows.

If $\alpha, \beta, \gamma \ldots$ are the roots of an algebraic equation, then the left side of (1) is the quotient of the $m$th and the $(m-1)$th power sum, and we thus obtain the theorem:

\begin{center}
\gesperrt{The quotient of the $m$th and the $(m-1)$th power sum approaches, as $m$ increases, the absolutely greatest among the roots.}
\end{center}

If one takes $m$ negative, and supposes $\alpha$ absolutely \gesperrt{smaller} than $\beta, \gamma \ldots$, it follows in the same way:

\begin{center}
\gesperrt{The quotient of the $-m$th and the $-(m+1)$th power sum approaches, as $m$ increases, the absolutely smallest among the roots.}
\end{center}

Since the power sums can be computed as symmetric functions from the coefficients, one need only take a sufficiently large $m$ in order to obtain an approximate value of the absolutely greatest and the absolutely smallest root.

But if among the roots of the equation there occur some which have the same absolute value, and this is the greatest or the smallest, then this method is not applicable — thus, for example, not to real equations in which a pair of imaginary roots occurs that surpasses the real ones in absolute value.

But if the greatest root, though it occurs singly, does not much surpass the next following one, then the procedure will yield some accuracy only slowly.

\gesperrt{Jacobi} has supplemented the \gesperrt{Bernoulli} method in one direction\footnote{Observatiunculae etc. Werke, Bd. 3, S. 280.}.
